\section{The local virtual degree}\label{sec:degree}

The putative relation of Lemma~\ref{lem:relation} would kill a
regular point orientation in the presence of other, possibly singular
and noncompact, components of the moduli space. We use an implicit atlas
on a neighborhood of the compact support of such a vanishing. First
we construct its pairing with compactly supported cohomology. We then
obtain the charts, including their orientation at nodal maps, from
fixed-target gluing. In this section references to numbered axioms,
deformation spaces and gluing charts are to \cite{Atlas}.

\subsection{A functional on a finite union of regular footprints}
In an implicit atlas, $X_I$ denotes the thickening belonging to a
finite set $I$ of indices, with obstruction space
$E_I=\bigoplus_{\alpha\in I}E_\alpha$, section $s_I:X_I\to E_I$,
and group $\Gamma_I=\prod_{\alpha\in I}\Gamma_\alpha$.
The footprint is the image of the zero locus under the quotient map
to the moduli space. The regular footprint uses the regular part of
the thickening. The local orientation system is the virtual
orientation system of the atlas.

\begin{proposition}\label{prop:functional}
Let $M$ be locally compact Hausdorff. Suppose we have the chart,
footprint and transversality data of an oriented topological implicit
atlas of virtual dimension $n$ on a neighborhood of a compact set.
We use no boundary on these charts. Let $O$ be an open set covered
by regular footprints of a finite set of indices $A$, allowing
multi-indices formed from $A$. All coefficients are rational.
There is a functional

\[
                   H_c^n(O)\longrightarrow\mathbb Q
\]

for this choice of charts. Suppose a point $p\in O$ has an oriented
manifold neighborhood of dimension $n$ in the zero locus. For some
index set $I\subset A$ assume its zero-coordinate lifts near $p$
are regular, are freely permuted in fibers of size $|\Gamma_I|$,
and that $s_I$ on this regular thickening is a submersion there.
Orient the neighborhood of $p$ by the virtual orientation. The
functional sends the point orientation class extended by zero to
$+1$.
\end{proposition}
\begin{proof}
For $I\subset J\subset A$ take the deformation
space $Y_{IJ}=X_{I,J,A}$ of Pardon's Definition 4.2.1. It is a
manifold with corners of dimension $n+\dim E_A+|I|$. Write
$e:Y_{IJ}\to E_A$ for its obstruction projection. On $e^{-1}(0)$
there is a continuous footprint map to the unthickened space.
Delete the part of $e^{-1}(0)$ lying over the complement of $O$,
an ordinary closed subset, and call the resulting open manifold
$Y^O_{IJ}$. Write $Z_{IJ}$ for its full zero locus and
$\pi:Z_{IJ}\to O$ for the footprint map. These restrictions commute
with the maps (4.2.1) for the one fixed set $A$. Away from $e=0$
no map to $O$ is required.

Let $O^+$ be the one-point compactification (with a disjoint point
if $O$ is compact). To a compact $K\subset O^+$ associate the
complex of relative singular chains

\[
 F_{IJ}^\bullet(K)=
 C_{n+\dim E_A-\bullet}
 \big(Y^O_{IJ},Y^O_{IJ}\setminus\pi^{-1}(K\cap O);
                     o_{E_A}^{\vee}\big)^{\Gamma_A}.
\]

Here $\pi^{-1}(K\cap O)$ is closed in $Y^O_{IJ}$, including when
$K$ contains infinity. Use the normalized transition maps of
Definition 4.2.6 and take their finite homotopy colimit $F^\bullet$.
This is the one-point compactification argument of Pardon's proof
of Lemma A.6.4 applied to each of the deformation spaces.

We give the descent checks. The assignments on compact subsets
are homotopy K-sheaves. The value at the empty set is zero, and the
two-compact-set axiom is singular-chain excision for the open
complements. For decreasing compact neighborhoods, a singular chain
has compact image. Its intersection with $Z_{IJ}$ is compact, and
the image of that intersection in $O^+$ is compact. If the chain
misses $\pi^{-1}(K\cap O)$, its footprint therefore misses a compact
neighborhood of $K$ as well. This proves continuity under shrinking
neighborhoods, including at infinity. Finite group invariants are
exact over $\mathbb Q$. Finite homotopy colimits preserve these
conditions.

At each $x\in O$, the stalk calculation is the one in Lemma 4.3.1
and Proposition 4.3.3. Excision away from the deleted zero loci
does not change the relative chains at $x$. The footprint diagrams
give $o_{M,x}$ in degree zero and zero in every other degree.
Coverage is required only over $O$. At infinity every relative
chain complex is zero. For the uniform lower bound in the purity
criterion, use Poincaré--Lefschetz duality with closed support in
the finite-dimensional manifolds $Y^O_{IJ}$, as in Lemma A.6.4.
The relative homology vanishes above their dimensions. There are
finitely many pairs and finitely many nondegenerate strings in
their indexing poset. Thus their shifted homotopy colimit still
has a uniform lower bound, also on neighborhoods of infinity.
The resulting degree-zero sheaf is $j_!o_O$ for $j:O\hookrightarrow
O^+$, by adjunction and the stalk calculations. Proposition A.5.4
now identifies $H^*(F^\bullet(O^+))$ with $H_c^*(O;o_O)$~\cite[\S\S4.2--4.3, Proposition A.5.4 and Lemma A.6.4]{Atlas}.

At $K=O^+$ the complement in each pair is the complement of the
\emph{whole} zero set in $Y^O_{IJ}$. Projection to $E_A$ therefore
defines a map to the relative chains of $(E_A,E_A\setminus0)$.
In the convention where the transition $J\to J'$ is the ordinary
pushforward divided by $|\Gamma_{J\setminus J'}|$, take the
projection on the $IJ$ term divided by $|\Gamma_J|$. Indeed
$|\Gamma_J|=|\Gamma_{J'}||\Gamma_{J\setminus J'}|$, and the
obstruction projection is unchanged by the transition. This gives
a chain map from the homotopy colimit by the augmentation on
strings of length zero, with value zero on higher strings. In
degree $n$, identify the resulting obstruction-space orientation
class with $1$. This gives the functional. Applying the projection
directly to chains relative to a proper compact subset of a larger
zero set would lack this map of pairs.

\par\smallskip\noindent\emph{Normalization at the regular point.} Let $V_I\subset O$ be the regular
footprint for $I$, intersected with $O$. In the diagonal term
$I=J$, take the slice $t_i=1$ for $i\in I$. It is the open part
of $E_{A\setminus I}\times X_I^{\rm reg}$ obtained by deleting
zeros over the complement of $O$. Its obstruction projection is
$(e',x)\mapsto(e',s_I(x))$. Its entire zero set is closed in
this open manifold and has proper finite quotient onto $V_I$.
Indeed it is the restriction of the finite group quotient of the
zero-coordinate chart to its invariant open regular locus.

Here is the slice calculation directly from Definition 4.2.1.
Denote the slice by $D$ and its zero set by $Z_D$. Before removal
of the zeros outside $O$ the closed set over a compact $K\subset O$
in this slice is

\[
       \{0\}\times\big((s_I|_{X_I^{\rm reg}})^{-1}(0)
                       \cap\psi_{\emptyset I}^{-1}(K)\big)
       \ \subset\ E_{A\setminus I}\times X_I^{\rm reg}.
\]

The pair uses the complement of this entire set. This detail
matters in using (4.3.14), whose displayed complement takes the
product of $E_{A\setminus I}$ with the complement in the second
factor. For example, take a point with trivial groups,
$X_\emptyset=\{p\}$ regular, and one additional chart
$X_\alpha=\mathbb R$ with $s_\alpha=\mathrm{id}$. For
$I=J=\emptyset$, Definition 4.2.5 gives
$H_1(\mathbb R,\mathbb R\setminus\{0\})=\mathbb Q$.
The pair $(\mathbb R,\varnothing)$ would give zero in degree one.

For any compact $K\subset O^+$ put $L=\pi^{-1}(K\cap O)$ in
$Y=Y_{II}^O$. On the interior $Y^\circ$, all $t_i>0$ and
$(t,e',x)$ are product coordinates, with $e_i=s_i(x)/t_i$.
In these coordinates the pair $(Y^\circ,Y^\circ\setminus L)$
is $(\mathbb R_{>0})^{I}$ times
$(D,D\setminus(L\cap D))$. Retraction of the first factor to
$(1,\ldots,1)$ proves the slice isomorphism on the interior.
The inclusion of this interior pair into $(Y,Y\setminus L)$
induces an isomorphism as well. To see it, apply Poincaré--Lefschetz
duality (A.6.6) to the closed set $L$ in $Y$ and to its intersection
with $Y^\circ$ in the open manifold $Y^\circ$. The coefficient
sheaf on $L$ in the first duality is the extension by zero of the
orientation sheaf from $L\cap Y^\circ$. Inclusion of chains is the
corresponding extension on compactly supported cohomology. This
naturality follows from the relative-chain map in the proof of
Lemma A.6.4, which is the orientation identification on interior
stalks and has zero source at boundary stalks. Both dual groups are
therefore the same. Tensoring with $o_{E_A}^\vee$ and taking
finite-group invariants preserves the isomorphism. The argument
also covers $K=O^+$, where $L$ is the whole zero set.

Applying Lemma A.6.4 once more to $Z_D$ in the slice, then taking
finite-group invariants as in (4.3.7), identifies the diagonal
term on $O^+$ with $H_c^*(V_I;o_O)$. The map of this term into
the homotopy colimit induces the identity of orientation sheaves
on $V_I$ by Proposition 4.3.3. Its source sheaf is extension by
zero from $V_I$. By adjunction this specifies the map of sheaves
on all of $O^+$, so on cohomology it is precisely extension by
zero $H_c^*(V_I;o_O)\to H_c^*(O;o_O)$.

Take a ball at $p$ on which the zero-locus cover splits into
$|\Gamma_I|$ sheets. Pullback of its point orientation to the
cover gives coefficient one on each sheet. In the relative chains
of the slice these classes are small oriented normal disks at the
lifts. Submersivity lets us choose them with boundary disjoint
from the entire zero set and obstruction projection of local degree
one. In chains with obstruction orientation coefficients, first
average one disk over $\Gamma_{A\setminus I}$, which fixes its
lift and acts simultaneously on the complementary obstruction
coordinates and their orientation coefficients. Its projected
degree remains one. Sum its translates over $\Gamma_I$ to obtain
the invariant representative on the free sheets. The
previous paragraph identifies its image in the full complex with
the point orientation extended by zero. Projection of the orbit sum
and division by $|\Gamma_I|$ therefore give one. This also works when regularity of
the unperturbed \emph{fixed-domain} operator fails. It uses an ordinary
submersion in one nonempty chart over the smooth ball.
\end{proof}
If the local orientation at $p$ vanished in $H_c^n(M)$, it would
already vanish with support in a compact set $K_0$ containing
$p$. Choose finitely many regular footprints covering $K_0$ and
an open $O$ inside their union containing $K_0$. Excision for
support in $K_0$ would make its image in $H_c^n(O)$ zero. The
functional contradicts this.

\subsection{An implicit atlas for maps in the fibers}
Let $\pi:E\to W$ be the smooth family of Lemma~\ref{lem:relation},
with $W$ open in $\CC^r$ and $n=2r$. Choose a smooth horizontal distribution
and give $E$ the almost-complex structure given by the vertical
structure and the lift of the standard structure on $W$.
Every connected compact holomorphic map to $E$ projects to a
constant. Holomorphic functions on each normalized component are
constant and agree at nodes. Thus unperturbed fixed-target maps
are exactly the maps into the fibers, including their Gromov
topology. Their Chern number in $E$ equals their vertical Chern
number since $TW$ is trivial.

The fixed-target construction can be used separately in each
homology class of maps to $E$. This class is locally constant in
the labeled-domain topology: on nearby domains collapse the
plumbing necks to the nodes, compare maps into a small tubular
neighborhood of the limiting graph, and use short geodesics in
$E$ to homotope to the limiting map composed with the collapse.
The collapse sends the fundamental cycle to the sum of the
fundamental cycles of the limiting components. Equivalently one
can integrate the corresponding integral currents in a compact
neighborhood and use local constancy of the integral current
homology class in the flat topology. A compact set of maps meets
only finitely many such sectors. The test also restricts to
nonconstant maps. A constant map has an image neighborhood
in a coordinate ball with an exact taming form, which excludes a closed
nonconstant holomorphic curve in that ball.

Use the indices of \cite[\S9.2]{Atlas} for the fixed target $E$ with \emph{vertical}
perturbation sections. Use actual codimension-two incidence disks
in $E$ and the usual marked-domain orbifold charts. The positivity
of rank of the differential on every unstable nonconstant component
gives the stabilizing disks of Lemma 9.2.8. In a small target
trivialization one can thicken a transverse disk in the fiber in
all base directions and round its edge away from the marked
intersections. For each finite index set $I$, take the fixed-target
thickening with these vertical perturbations. Every solution, at
arbitrary vertical obstruction coordinates, still projects
holomorphically to $W$ and lies in a single fiber. The spaces, group
actions, footprint quotients, obstruction coordinates and forgetful
maps are therefore the ordinary data of Definitions 9.2.1--9.2.4.
The finite group need not act effectively on its marked-domain chart.

Choose indices on sufficiently small marked-domain charts. On each
one we can also specify a complex $\Gamma_\alpha$-representation
$P_\alpha$ of dimension $kr$, where $k$ is the fixed arithmetic
genus, and equivariant perturbation forms whose projection induces
an isomorphism

\[
           P_\alpha \ \simeq\
             H^1(T,\mathcal O_T)\otimes\mathbb C^r
\]

on each labeled domain $T$ in the chart. Here is the construction.
At the central stabilized curve use the right-hand space as the
representation. Smooth $(0,1)$ forms supported in the smooth locus
away from the marked points span its first factor by duality with
dualizing differentials. Such a differential vanishing on the smooth
locus is zero, including at every branch of a node. Choose a complex
linear right inverse and average over the finite group. The forms
extend on disks in the smooth locus of the universal curve. Extend
and average again, keeping the original equivariant values at the
center. The classes still give a basis after shrinking the chart.
Indeed the dual Hodge bundle is locally free for the proper flat
family of connected nodal curves of fixed genus. In a Hodge frame
the pairing with these disk forms is continuous up to the central
fiber, by integration over a fixed compact set away from the nodes.
The extension may be chosen smooth on the total smooth locus.
This also handles normalization components on which the map is
constant. Take the horizontal lifts of the resulting
$\mathbb C^r$-valued forms on $\mathcal C_\alpha\times E$. They
are allowed perturbations in Definition~9.2.1(vi) of \cite{Atlas}. For relative duality and base change see \cite[Tags 0E6N and 0E6S]{Stacks}.

For example, take the rational nodal genus-one curve obtained by
identifying $0$ and $\infty$, with a stabilizing mark at $1$. The
involution $\tau(z)=1/z$ acts by $-1$ on the dualizing differential
$dz/z$, hence by $-1$ on $H^1(\mathcal O_T)$. Choose a form
$\phi$ supported on a disk near $2$ pairing nontrivially with
$dz/z$. The form $(\phi-\tau^*\phi)/2$ is anti-invariant and has
the same nonzero pairing. An invariant horizontal coordinate
would miss this cokernel. Taking $P_\alpha$ with its actual
stabilizer representation deals with this node without passing to
an unlabeled quotient.

For each nonempty $I$ select for the moment $\alpha\in I$. Replace
the datum at $\alpha$ by the same datum with obstruction space
$E_\alpha\oplus P_\alpha$ and the additional horizontal perturbations.
Keep every other datum unchanged. This is an ordinary collection
of Pardon's fixed-target indices. On any simultaneous thickening
containing $\alpha$ the projected perturbed equation reads

\[
        \bar\partial w+\eta_\alpha(h)=0,\qquad
        w=\pi u,\quad h\in P_\alpha .
\]

Taking its Dolbeault class gives $h=0$, and then $w$ is constant.
Conversely every solution of the original vertical thickening
is a solution with $h=0$. This is a homeomorphism of the two
thickenings. All marks and domain labels are identical, the graph
topologies agree, and the extra coordinate on one side is constantly
zero. If the distinguished index is retained, the homeomorphism
also intertwines the footprint maps on the locus where the new
obstruction coordinates are zero.

The regularity test in Definition 9.2.5 is independent of which
$\alpha\in I$ is chosen for this purpose. At a solution the
horizontal block of the fixed-domain linearization is
$(\xi,h)\mapsto\bar\partial\xi+\eta_\alpha(h)$. Its cokernel is
zero, and its kernel is the constant horizontal displacement
$T_bW$. Derivatives of the horizontal forms with respect to domain
labels have coefficient $h=0$. In the kernel of this block the
remaining equation is exactly the vertical linearization with
free base value $b$. Transverse intersections with incidence disks,
used to define the labels in (9.2.3), move smoothly with the map
on the fixed domain. Thus the full augmented operator is surjective
if and only if the vertical operator with free base is surjective.
Triviality of the automorphism group is the same on both sides.

Define regularity on nonempty thickenings by this intrinsic test,
and set the regular locus of the empty chart to the empty set.
For $I\subset J$ with $I$ nonempty, augment \emph{one index in $I$}
on both sides of the inclusion. Under the two homeomorphisms just
proved, fixed-target Proposition 9.2.6 gives the required openness
and the topological submersion onto $E_{J\setminus I}$. The
auxiliary space never occurs in this quotient. For $I$ empty the
submersion and regular-footprint freeness requirements have empty
domains of verification. The footprint and quotient axioms still
hold there without a regularity assumption. The dimension in the
augmented model is

\[
       2r(1-k)+\dim_{\mathbb R} E_I+2rk=n+\dim_{\mathbb R}E_I
\]

for unmarked maps in the Chern-zero sector.

For coverage choose a stabilizing disk and finite-dimensional
vertical perturbations spanning the fixed-domain \emph{vertical}
cokernel, as in the proof of Lemma 9.2.9. At the unperturbed map
the label derivatives of these vertical perturbations have zero
coefficients. Perturbations supported away from all special points
suffice by that lemma's elliptic argument. Extend the sections
locally in the vertical tangent bundle and use the induced finite
group representation for equivariance. The choice of a labeled
domain chart removes the residual automorphisms of the lifted object.
This produces a nonempty regular footprint at every point. The construction
uses compactness of the domain curves, and only finitely many indices
are needed near a compact $K_0$. Pardon's compact-target assumption
in the global GW application supplies compactness of the entire
unperturbed space. Proposition 9.2.6 itself is a gluing calculation
near a single compact map. Its stabilized fixed-target model in
Appendix B.1 uses a smooth almost-complex target and transverse
disks. All the maps in this local calculation stay near the fixed compact image. Local compactness of the unperturbed
space also follows from the regular footprints themselves. Each is
an open set in the Hausdorff moduli space and the finite-group
quotient of the closed zero locus in a locally compact regular
chart. These footprints cover.

\subsection{Orientation on overlaps} At a smooth-domain regular
point of $X_I$, retain the same domain labels and incidence disks.
For $\alpha,\beta\in I$, write $H_T=H^1(T,\mathcal O_T)\otimes
\mathbb C^r$ and denote the two maps to this quotient by
$A_\alpha:P_\alpha\to H_T$ and $A_\beta:P_\beta\to H_T$.
Identify $P_\alpha$ with $P_\beta$ by the complex isomorphism
$G=A_\beta^{-1}A_\alpha$. On the common linearized domain consider
the horizontal perturbation forms

\[
            \eta_s=(1-s)\eta_\alpha+s\,\eta_\beta G,
            \qquad 0\leq s\leq1 .
\]

Every $\eta_s$ projects to $A_\alpha$. Thus
$(\xi,h)\mapsto\bar\partial\xi+\eta_s(h)$ is onto for the whole
path, and its kernel always has $h=0$ and $\xi$ constant. This
also describes the horizontal equation with domain variations,
since the projection of the original map is constant and the
horizontal perturbation coefficient at the solution is zero.
On this kernel the remaining linear equation, including the old
perturbations and moving incidence marks, is the same for every
$s$. Regularity therefore gives a path of surjective linearizations
with the very same tangent kernel. The canonical determinant
orientation varies continuously along that path in the fixed
kernel. Identifying its two auxiliary factors by $G$ preserves
their real orientations. This proves equality of the two
orientations of the regular thickening at the smooth point.
Only a path of linear operators at that point is needed. The
construction commutes with relabeling, since $G$ is defined by
the induced maps to the same horizontal cohomology group.

Each augmented chart already has a continuous orientation
identification by Proposition 9.2.6(iii). Both live on the same
topological regular locus by the homeomorphisms above. By \cite[\S B.12]{Atlas}
the plumbing parameters can all be nonzero arbitrarily close to
every nodal point of that locus. The two identifications therefore
agree there as well, since their ratio is a locally constant sign.
For $I\subset J$ with $I$ nonempty, choose an index in $I$ on both
sides. On the regular footprint the forgetful map at zero extra
coordinates is a finite covering of an open set in $X_I^{\rm reg}$.
Smooth-domain lifts are dense on this covering. On these lifts the
ordinary determinant exact sequence for adding $E_{J\setminus I}$
gives the orientation transition for the vertical charts, since
the same oriented $P_\alpha$ cancels from both sides. Continuity
gives the transition at the nodal lifts. For a triple of nested
nonempty index sets one may choose the same index in the smallest
one throughout, so the transition maps compose. All extra complex
factors have even real dimension. This verifies orientation descent
on the regular nodal overlaps used by the compact-support test~\cite[Proposition 9.2.6 and \S\S B.12--B.13]{Atlas}.

\begin{proof}[Proof of Theorem~\ref{thm:independence}]
Use the regular point over the least genus in Lemma~\ref{lem:relation}.
Its normal bundle is acyclic, its map automorphism group is trivial,
and the family is constant over a ball at its marked parameter.
In any nonempty regular chart covering this point the map to the
vertical obstruction coordinates is a submersion near its lift.
Indeed the unperturbed linearization with domain variations is onto.
On this smooth domain the implicit function theorem solves the perturbed
equation with arbitrary prescribed small obstruction coordinates.
Transverse intersections recover the stabilizing marks. The kernel
is the tangent to the graph of the unperturbed solutions, oriented
by the complex normal operator and the base. There are $|\Gamma_I|$
lifts. Proposition~\ref{prop:functional} therefore pairs the local
point orientation to $+1$ on a neighborhood of every compact set
containing it. The vanishing of Lemma~\ref{lem:relation} is impossible.
\end{proof}
